% GENERATED FILE. Edit canonical lessons, then run npm run generate:books.
% canonical-source-hash: fnv1a64:4c72f050c9241f24
% canonical-lessons: SA-W65-letter-nga, SA-R65-letters-first-pass, SA-R65-letters-second-pass

\chapter{The Last Letter You Read}
\label{ch:sa-last-letter-read}
\hlchaptermodality{65}

\begin{chapteropening}
\textbf{By the end of this chapter:} \emph{I can write the letter nga, and find it in the word for number.}

Complete the last lesson of chapter 65: i can write the letter nga, and find it in the word for number.
\end{chapteropening}

\section[ṅa]{\sk{ङ} (ṅa) — the nasal before \sk{ख} in \sk{सङ्ख्या}}

\label{lesson:SA-W65-letter-nga}



\begin{quote}
Before the new one: write \sk{इ} and \sk{ऋ} once each, and name them.

Now say \textbf{\sk{सङ्ख्या}} (\emph{saṅkhyā}), \textquotedblleft{}number\textquotedblright{}. You have read it for a long time. This lesson takes one letter out of it and puts it in your hand.
\end{quote}

\subsection*{Script you'll notice: \sk{ङ}}
\textbf{\sk{ङ}} — \emph{ṅa}, the nasal made where \emph{k} and \emph{kh} are made.

It is inside \textbf{\sk{सङ्ख्या}} (\emph{saṅkhyā}), \textquotedblleft{}number\textquotedblright{}, from the chapter on numbers. Say \emph{saṅkhyā}: the \emph{ṅ} sits at the back of the mouth, ready for the \emph{kh} after it. Sanskrit spells that nasal out as \textbf{\sk{ङ}}; Hindi writes the same word \textbf{\sk{संख्या}}, with a dot above instead.

\subsection*{Writing: \sk{ङ}}
Put your pen on \sk{ङ} and follow its line. Copy the shape you can see — slowly, and larger than it is printed.

\begin{quote}
This book does not yet tell you \textbf{where to start this letter or which way to travel}. It is not written down here until it can be written down with a source. Copying what is in front of you needs no such source, and it is how the shape gets into your hand in the meantime.
\end{quote}

\subsection*{Your turn}
\begin{itemize}
\raggedright
  \item \emph{Look:} at \textbf{\sk{सङ्ख्या}} and point at \sk{ङ} inside it
  \item \emph{Trace it:} \sk{ङ} three times, saying \emph{ṅ} as you finish each one
  \item \emph{Write it:} \sk{ऋ}, then \sk{ङ}, from memory
  \item \emph{Read it:} \textbf{\sk{सङ्ख्या}} aloud once more
\end{itemize}

\begin{tcolorbox}[breakable,colback=teal!4,colframe=teal!35!black,title={Before you move on}]
Which letter is this — \sk{ङ}? (\textbf{ṅa}, \emph{ṅ}.) Which word did it come from? (\textbf{\sk{सङ्ख्या}}, \emph{saṅkhyā}, \textquotedblleft{}number\textquotedblright{}.)
\end{tcolorbox}
\section[ekam akṣaram]{(one letter) — From the letter to the word — a first pass}

\label{lesson:SA-R65-letters-first-pass}



\begin{quote}
Say \textbf{\sk{सङ्ख्या}} (\emph{saṅkhyā}), then write the letter it gave you.
\end{quote}

\subsection*{Your turn}
\noindent
\begin{tabularx}{\linewidth}{@{}>{\raggedright\arraybackslash}X>{\raggedright\arraybackslash}X>{\raggedright\arraybackslash}X@{}}
\toprule
\textbf{letter} & \textbf{name} & \textbf{the word it came from} \\
\midrule
\sk{ङ} & \emph{ṅa} & \textbf{\sk{सङ्ख्या}} \emph{saṅkhyā} \\
\bottomrule
\end{tabularx}

\begin{itemize}
\raggedright
  \item \emph{Write it:} each letter in the table from memory, then check it
  \item \emph{Say it:} the word each one came from
  \item \emph{Read it:} each word in the table aloud, and point at its letter
\end{itemize}

\begin{tcolorbox}[breakable,colback=teal!4,colframe=teal!35!black,title={Before you move on}]
Which word did each letter come from? (\textbf{\sk{ङ}} from \textbf{\sk{सङ्ख्या}}.)
\end{tcolorbox}
\section[ekam akṣaram]{(one letter) — From the word back to the letter — a second pass}

\label{lesson:SA-R65-letters-second-pass}



\begin{quote}
Write \textbf{\sk{ङ}} without looking, then say the word it came from.
\end{quote}

\subsection*{Your turn}
\noindent
\begin{tabularx}{\linewidth}{@{}>{\raggedright\arraybackslash}X>{\raggedright\arraybackslash}X>{\raggedright\arraybackslash}X@{}}
\toprule
\textbf{letter} & \textbf{name} & \textbf{the word it came from} \\
\midrule
\sk{ङ} & \emph{ṅa} & \textbf{\sk{सङ्ख्या}} \emph{saṅkhyā} \\
\bottomrule
\end{tabularx}

\begin{itemize}
\raggedright
  \item \emph{Say it:} each word in the table, then write the letter it gave you without looking
  \item \emph{Write it:} each letter once more, slowly, larger than it is printed
\end{itemize}

\begin{tcolorbox}[breakable,colback=teal!4,colframe=teal!35!black,title={Before you move on}]
Say each word in the table once more, and name the letter you took from it.
\end{tcolorbox}
